


\section{Basic Terminology and Notations}
\label{sec:6.1}

A \textsl{state vector\/} for a given set is the indicator of that set, i.e., it is one at the indices which belong to the given set, and zero
otherwise. We denote by $\chi_U$ the state vector for a set $U$. The vector $u$ is $1\times n$, all entries equal to 1.
$J=\dg  u.u$ is the all-ones matrix.

If $\A=R+B$ is the adjacency matrix of the graph, $A=\frac12\,\A$ is stochastic with invariant distribution $\pi$. 
The corresponding Cesaro limit matrix is $\Omega=\dg u.\pi$. The diagonal matrix $\omega$ has diagonal entries
$\omega_{ii}=\pi_i$, the invariant distribution. Set $\Delta = \frac12\,(R-B)$.

The semigroup $\cS$ is generated by $R$ and $B$. The kernel of $\cS$ is $\cK$. The idempotents of $\cK$ are $\E(\cK)$.
For $k\in\cK$, the corresponding \textsl{range vector\/} is $\rho(k)=\chi_{{\rm range\,}(k)}$, the state vector for the range of $k$.
With $r$ the rank of $\cK$, $\rho(k)\dg u.=r$ for every $k\in\cK$.


We denote by $\alpha$ the measure on the idempotents induced from the partitions of $\cK$. 
Denote by $\beta$ the measure on the ranges of the kernel $\cK$.

We use the summation convention wherein Greek indices are summed over regardless of position.
